
\(K^\bullet(T)\) et la
filtration par le degré sur le deuxième membre résulte aisément de (V 6.6.4) et du fait que les \(c^i\) définissent une \emph{théorie de Chern} de \(K^\bullet(T)\) au sens de (V 6.13), lorsqu'on considère \(K^\bullet(T)\) comme une \(\lambda\)-algèbre augmentée sur l'anneau binômial \(H^0(T,\mathbb Z)\)). On en déduit des applications
\begin{equation}\tag{7.13.7}
  \operatorname{c\ell}:\Gr^i(T)\longrightarrow H^{2i}(T,\mathbb Q_\ell(i))
  \stackrel{\mathrm{dfn}}{=}
  H^{2i}(T,\mathbb Z_\ell(i))\otimes_{\mathbb Z_\ell}\mathbb Q_\ell,
\end{equation}
par la formule
\begin{equation}\tag{7.13.8}
  \operatorname{c\ell}(x^{(i)})=
  (1/(-1)^{i-1}(i-1)!)\,\psi_i(x^{(i)}),
\end{equation}
comparer (XIV 5.1 et 5.2). Ceci posé, la commutativité des diagrammes (7.13.2) et (7.13.4) implique trivialement celle des diagrammes analogues
\begin{equation}\tag{7.13.9}
\begin{tikzcd}[column sep=large,row sep=large]
\Gr^i(X_t)_{\mathbb Q} \arrow[r,"\operatorname{c\ell}"] \arrow[d,"\sigma_{\mathbb Q}"'] &
H^{2i}(X_t,\mathbb Q_\ell(i))\\
\Gr^i(X_s)_{\mathbb Q} \arrow[r,"\operatorname{c\ell}"] &
H^{2i}(X_s,\mathbb Q_\ell(i)) \arrow[u,"\sigma^{-1}"]
\end{tikzcd}
\;,
\end{equation}
\begin{equation}\tag{7.13.10}
\begin{tikzcd}[column sep=large,row sep=large]
\Gr^i(X_{\bar t})_{\mathbb Q} \arrow[r,"\operatorname{c\ell}"] \arrow[d] &
H^{2i}(X_{\bar t},\mathbb Q_\ell(i)) \arrow[d,"\wr"]\\
\Gr^i(X_{\bar s})_{\mathbb Q} \arrow[r,"\operatorname{c\ell}"] &
H^{2i}(X_{\bar s},\mathbb Q_\ell(i))
\end{tikzcd}
\;.
\end{equation}
\textbf{Remarque 7.13.11.} La commutativité des deux derniers diagrammes, qu'on peut exprimer en disant que la spécialisation des cycles algébriques est compatible avec la formation des classes de cohomologie \(\ell\)-adiques associées, n'est pas très satisfaisante sous la forme énoncée, puisqu'on a été obligé de remplacer les coefficients \(\mathbb Z_\ell(i)\) par \(\mathbb Q_\ell(i)\), c'est-à-dire de négliger la torsion. Pour voir un énoncé qui ne néglige pas la torsion, il semble qu'on
